\SourceLeaf{89}
\Article{1}
Ad explicationem phaenomenorum magneticorum duo fluida magnetica postulamus; alterum cum physicis vocamus boreale, alterum australe. Elementa fluidi alterius attrahere alterius elementa, contra bina elementa eiusdem fluidi mutuo se repellere supponimus, et quidem utramque actionem variari in ratione inversa quadrati distantiae. Veritatem huius legis per ipsas nostras observationes plenissime confirmari infra apparebit.

Fluida ista non per se apparent, sed tantummodo iuncta cum particulis ponderabilibus corporum talium, quae magnetismi capaces sunt, illorumque actiones in eo se manifestant, quod has vel ad motum sollicitant, vel motum, quem aliae vires in ipsas agentes, e. g. gravitas, producerent, impediunt vel mutant.

Actio itaque quantitatis datae fluidi magnetici in quantitatem datam vel eiusdem fluidi vel alterius in distantia data comparabilis erit cum vi motrice data, i. e. cum actione vis acceleratricis datae in massam datam, et quum fluida magnetica ipsa non nisi per effectus quos producunt cognoscere liceat, hi ipsi illorum mensurae inservire debent.

Quo igitur hanc mensuram ad notiones distinctas revocare possimus, ante omnia circa tria quantitatum genera unitates stabilire oportet, puta unitatem distantiarum, unitatem massarum ponderabilium, unitatem virium acceleratricium. Pro tertia accipi potest gravitas in loco observationis: quod si minus arridet, insuper accedere debet unitas temporis, eritque nobis vis acceleratrix ea $=1$, quae in unitate temporis mutationem velocitatis corporis in ipsius directione moti unitati aequalem gignit.

\SourceLeaf{90}
His ita intellectis, unitas quantitatis fluidi borealis ea erit, cuius vis repulsiva in aliam ipsi aequalem in distantia $=1$ positam aequivalet vi motrici $=1$, i. e. actioni vis acceleratricis $=1$ in massam $=1$, idemque de unitate quantitatis fluidi australis valebit: in hac determinatione manifesto tum fluidum agens, tum fluidum in quod agitur, in punctis physicis concentrata concipi debent. Insuper autem supponere oportet, attractionem inter quantitates datas fluidorum heteronymorum in distantia data aequalem esse repulsioni inter quantitates resp. aequales fluidorum homonymorum. Actio itaque quantitatis $m$ fluidi magnetici borealis in quantitatem $m'$ eiusdem fluidi in distantia $r$ (dum utrumque in puncto physico concentratum supponitur), exprimitur per $\frac{mm'}{rr}$, sive vi motrici $=\frac{mm'}{rr}$ in directione a priori versus posterius agenti aequivalet, manifestoque haec formula generaliter valet, si, quod semper abhinc subintelligemus, quantitas fluidi australis tamquam negativa spectatur, ubi valor negativus vis $\frac{mm'}{rr}$ pro repulsione attractionem indicabit.

Si itaque in puncto physico aequales quantitates fluidi borealis et australis simul adsunt, nulla omnino actio hinc orietur, si vero inaequales, excessus alterius tantum, quem magnetismum \textit{liberum} (positivum seu negativum) vocabimus, in considerationem veniet.

\Article{2}
Hisce suppositionibus fundamentalibus adhuc aliam, quam experientia undique confirmat, adiicere oportet, scilicet, quodvis corpus, in quo fluida magnetica adsint, semper aequalem utriusque quantitatem continere. Quinadeo experientia docet, hancce suppositionem etiam ad singulas talis corporis partes quantumvis parvas, dummodo sensibus nostris discerni possint, extendendam esse. Sed quum per ea, quae in fine art. praec. monuimus, actio eatenus tantum existere possit, quatenus aliqua separatio fluidorum locum habet, necessario hanc per intervalla tam parva fieri supponere debemus, ut mensuris nostris non sint accessibilia.

Corpus itaque magnetismi capax concipi debet tamquam compages innumerarum particularum, quarum quaevis certam quantitatem fluidi magnetici borealis et aequalem australis contineat, ita quidem, ut vel uniformiter inter se mixtae sint (magnetismus lateat), vel separationem minorem maioremve iniverint (magnetismus evolutus sit), quae tamen separatio numquam in transfusionem fluidi ab \SourceLeaf{91}una particula in aliam abire potest. Nihil refert, utrum separatio maior a maiori quantitate fluidorum quae libera evaserunt, an a maiori intervallo interposito orta supponatur: manifesto autem propter magnitudinem separationis simul eius directio in considerationem venire debet, quae prout in diversis corporis particulis vel conspirat vel refragatur, maior minorve energia totalis respectu punctorum extra corpus oriri poterit.

Quomodocunque autem distributio magnetismi liberi intra corpus se habeat, semper eius loco, per theorema generalius, substituere licet secundum certam legem aliam distributionem in sola corporis superficie, quae respectu virium extrorsum agentium illi exacte aequivaleat, ita ut elementum fluidi magnetici extra corpus ubicunque positum prorsus eandem attractionem vel repulsionem experiatur a distributione magnetismi vera intra corpus atque a fictitia in eius superficie. Eandem fictionem ad \textit{bina} corpora, quae ratione magnetismi liberi in ipsis evoluti in se invicem agunt, extendere licet, ita ut pro utroque distributio fictitia in superficie distributionis verae internae vice fungi possit. Hocce demum modo vulgari loquendi mori, qui e. g. alteri acus magneticae extremitati solum magnetismum borealem, alteri australem tribuit, sensum verum conciliare possumus, quum manifesto haec phrasis cum principio fundamentali supra enunciato, quod alia phaenomena imperiose postulant, non quadret. Sed haec obiter hic annotavisse sufficiat; de theoremate ipso, quum ad institutum praesens non sit necessarium, alia occasione copiosius agemus.

\Article{3}
\textit{Status magneticus} corporis consistit in ratione distributionis magnetismi liberi in singulis eius particulis. Respectu mutabilitatis huius status discrimen essentiale inter corpora diversa magnetismi capacia animadvertimus. In aliis, e. g. in ferro molli, ille status per levissimam vim protinus mutatur, hacque cessante status anterior redit: contra in aliis, praesertim in chalybe durato, vis certam intensitatem attigisse debet, antequam sensibilem status magnetici mutationem producere possit, vique cessante corpus vel in statu quem acquisivit permanet, vel saltem ad priorem non ex asse revenit. In corporibus itaque prioribus moleculae fluidi magnetici semper ad aequilibrium perfectum virium, quae tum inter ipsa mutuo, tum a caussis externis emanant, se componunt, vel saltem a tali aequilibrio sensibiliter vix differunt: contra in corporibus posterioris generis sta%
